%%═══════════════════════════════════════════════════════════════════
%% Lecture 07 — Two-Body Decays, Scattering Cross-Sections,
%%              Symmetry Factors, and Introduction to Fermions
%% 77609 — Introduction to Elementary Particles
%% Date: 30/04/2026
%% Lecturer: Prof. Yonit Hochberg
%%═══════════════════════════════════════════════════════════════════

\documentclass[../main.tex]{subfiles}
\usepackage{preamble}

\begin{document}

\section{Lecture 07 — Two-Body Decays, Cross-Sections, Symmetry Factors, and Fermions}\label{sec:lec07}
\begin{center}
\begin{tabular}{ll}
  \textbf{Date:}\quad 30/04/2026   & \\
\end{tabular}
\end{center}
\hrule\vspace{1.5em}
%%──────────────────────────────────────────────────────────────────
\subsection*{Overview}

This lecture completes the treatment of phase space and observables, then opens the chapter on fermions.

\begin{enumerate}
  \item \textbf{Two-body decay kinematics}: explicit formulae for the energies and three-momenta of the daughter particles in the rest frame; integration over the solid angle; the compact two-body decay rate formula.
  \item \textbf{General scattering cross-section}: the phase-space formula for $1+2 \to 3+\cdots+n$; the Lorentz-invariant prefactor; the standard form of $d\sigma$.
  \item \textbf{$2\to2$ scattering in the center-of-mass frame}: working out $s$, $t$, $u$ for $\phi^3$ theory; the equal-mass simplification; the full differential cross-section.
  \item \textbf{Symmetry factors via contraction counting}: why the three channels in $\phi\phi\to\phi\phi$ enter with equal weight; the concept of contractions.
  \item \textbf{Dimensional analysis of matrix elements}: $[\mathcal{M}] = \text{mass}^{4-n}$ for an $n$-particle process; sanity checks on $\DecayRate$ and $\sigma$.
  \item \textbf{Summary of perturbative QFT}: free vs.\ interacting fields; role of Feynman diagrams.
  \item \textbf{Introduction to fermions}: the Dirac equation; gamma matrices and Clifford algebra; $\gamma^5$; Dirac conjugate $\bar\psi$; the Dirac Lagrangian; Weyl spinors.
\end{enumerate}

%%──────────────────────────────────────────────────────────────────
\subsection*{Two-Body Decay Kinematics}

\noindent\textbf{Physical intuition:} When a particle at rest decays to exactly two daughters, four-momentum conservation in 3D leaves no freedom in the magnitude of the momenta — the daughters must fly back-to-back with equal and opposite three-momenta, completely determined by the three masses.  The only remaining freedom is the direction of the decay axis, which is integrated over to obtain the total rate.

\subsubsection*{Setup}

Consider particle~1 of mass $M_1$ at rest decaying to particles~2 and~3:
\[
  1 \;\xrightarrow{\text{rest}}\; 2 + 3.
\]
The four-momentum of particle~1 is $p_1 = (M_1,\,\bvec{0})$, since it is at rest.

\subsubsection*{Energies and momenta of the daughters}

Conservation of four-momentum, $p_1 = p_2 + p_3$, together with the on-shell conditions $p_i^2 = m_i^2$ uniquely determines:
\begin{equation}
  \boxed{E_2 = \frac{M_1^2 + m_2^2 - m_3^2}{2M_1},}
  \label{eq:lec07:E2}
\end{equation}
and (by the same calculation with $2 \leftrightarrow 3$)
\begin{equation}
  E_3 = \frac{M_1^2 + m_3^2 - m_2^2}{2M_1}.
  \label{eq:lec07:E3}
\end{equation}
Because three-momentum is conserved and the parent is at rest, $\bvec{p}_2 + \bvec{p}_3 = \bvec{0}$, so $|\bvec{p}_2| = |\bvec{p}_3| \equiv p_*$.  The magnitude is:
\begin{equation}
  \boxed{
  p_* = \frac{\sqrt{
    \bigl[M_1^2 - (m_2+m_3)^2\bigr]\bigl[M_1^2 - (m_2-m_3)^2\bigr]
  }}{2M_1}.
  }
  \label{eq:lec07:pstar}
\end{equation}
This can be written more compactly using the triangle (Källén) function $\lambda(a,b,c)=(a-b-c)^2-4bc$, but the form above makes the physics transparent: $p_* > 0$ requires $M_1 > m_2 + m_3$ (the decay threshold).

\subsubsection*{Integration over the solid angle}

From Lecture~06, the standard decay rate formula (with the $p_j^0$ integrals already done) gives the differential decay rate:
\begin{equation}
  d\DecayRate = \frac{1}{32\pi^2}\,|\Mamp|^2\,
    \frac{p_*}{M_1^2}\,d\Omega_2,
  \label{eq:lec07:dGamma2body}
\end{equation}
where $d\Omega_2 = \sin\theta_2\,d\theta_2\,d\phi_2$ is the solid angle of particle~2.  If $|\Mamp|^2$ is independent of the decay direction (true for scalar particles), we integrate:
\[
  \int d\Omega = \int_{-1}^{1}d(\cos\theta)\int_0^{2\pi}d\phi = 4\pi.
\]
The total two-body decay rate is then:
\begin{equation}
  \boxed{
  \DecayRate_{1\to 23} = \frac{p_*}{8\pi\,M_1^2}\,|\Mamp|^2.
  }
  \label{eq:lec07:Gamma2body}
\end{equation}

\begin{remark}[Kinematics fully determined]
For a two-body decay, all momenta are uniquely determined by the masses up to the overall direction.  In contrast, for three-body and higher decays there is a continuous distribution of final-state configurations — the Dalitz plot.  The two-body case is therefore both the simplest and the most predictive.
\end{remark}

%%──────────────────────────────────────────────────────────────────
\subsection*{General Scattering Cross-Section}

\noindent\textbf{Physical intuition:} The cross-section is the effective area that a target particle presents to an incoming beam.  It is defined so that the rate equals the flux times the cross-section.  The factor $1/(2E_1\,2E_2\,|\bvec{v}_1-\bvec{v}_2|)$ in the formula is exactly this flux factor, and the rest is the same phase-space integral as in the decay rate.

\subsubsection*{The general formula}

For the $2\to n$ scattering process $1+2 \to 3+\cdots+n$, the total cross-section is:
\begin{equation}
  \sigma = \frac{1}{2E_1\,2E_2\,|\bvec{v}_1-\bvec{v}_2|}
  \int |\Mamp|^2\,
  (2\pi)^4\,\delta^{(4)}\!\Bigl(p_1+p_2 - \sum_{j=3}^{n}p_j\Bigr)
  \prod_{j=3}^{n}\frac{d^3p_j}{(2\pi)^3\,2E_j},
  \label{eq:lec07:sigma-general}
\end{equation}
where $E_j = \sqrt{|\bvec{p}_j|^2+m_j^2}$ throughout.  The three kinematic constraints — on-shell, positive energy, and four-momentum conservation — are built into the measure, exactly as in the decay rate.

\subsubsection*{Origin of the prefactor}

The prefactor $1/(2E_1\,2E_2\,|\bvec{v}_1-\bvec{v}_2|)$ has two origins:
\begin{itemize}
  \item $1/(2E_1)\cdot 1/(2E_2)$: normalization of the two initial-state wave functions (same $\sqrt{2\omega_k}$ normalization used for the decay rate).
  \item $1/|\bvec{v}_1-\bvec{v}_2|$: the cross-section is obtained by dividing the transition rate by the flux of incoming particles; the flux is proportional to the relative velocity $|\bvec{v}_1-\bvec{v}_2|$.
\end{itemize}

\subsubsection*{Lorentz-invariant form of the prefactor}

While individual $E_i$ and $|\bvec{v}_i|$ are frame-dependent, the combination
\begin{equation}
  2E_1\,2E_2\,|\bvec{v}_1-\bvec{v}_2|
  = 4\sqrt{(p_1\cdot p_2)^2 - m_1^2\,m_2^2}
  \label{eq:lec07:flux-Lorentz}
\end{equation}
is Lorentz invariant (it is a Lorentz scalar).  This makes the total cross-section Lorentz invariant, as it must be: the total number of scattering events is a physical observable, independent of frame.

%%──────────────────────────────────────────────────────────────────
\subsection*{$2\to 2$ Scattering: Differential Cross-Section}

\subsubsection*{General result}

For $1+2\to3+4$, the differential cross-section with respect to the solid angle in the center-of-mass (CM) frame is:
\begin{equation}
  \dcross\bigg|_{\rm CM}
  = \frac{1}{2E_1\,2E_2\,|\bvec{v}_1-\bvec{v}_2|}
    \cdot \frac{|\bvec{p}_3|}{(2\pi)^2\,4\sqrt{s}}\,
    |\Mamp(1+2\to3+4)|^2,
  \label{eq:lec07:dsigma-22}
\end{equation}
where $\sqrt{s}$ is the CM energy.

\subsubsection*{Equal-mass simplification}

When all four particles have the same mass $m$ (including the massless limit $m\to0$), the CM momenta satisfy $|\bvec{p}_3| = |\bvec{p}_1| \equiv p_{\rm CM}$, and \eqref{eq:lec07:dsigma-22} simplifies to:
\begin{equation}
  \boxed{
  \dcross\bigg|_{\rm CM}
  = \frac{|\Mamp|^2}{64\pi^2\,s}.
  }
  \label{eq:lec07:dsigma-equal-mass}
\end{equation}
This is the form most frequently used in practice.

%%──────────────────────────────────────────────────────────────────
\subsection*{Worked Example: $\phi\phi\to\phi\phi$ in $\phi^3$ Theory}

\noindent\textbf{Physical intuition:} In $\phi^3$ theory, two scalar particles can scatter by exchanging a virtual $\phi$, and the intermediate virtual particle can carry different four-momenta depending on which pair of external legs it connects.  The three distinct topologies ($s$, $t$, $u$ channels) correspond to these three choices, and in the CM frame they differ by how the scattering angle $\theta$ enters.

\subsubsection*{Review of the three diagrams}

The interaction Lagrangian is $\lagr_{\rm int} = \lambda\phi^3$.  For $\phi\phi\to\phi\phi$ there are three Feynman diagrams at order $\lambda^2$, labeled by the Mandelstam variable of the propagator:
\begin{align}
  \Mamp_s &\sim \frac{\lambda^2}{s - m^2}, & s &= (p_1+p_2)^2, \label{eq:lec07:Ms}\\
  \Mamp_t &\sim \frac{\lambda^2}{t - m^2}, & t &= (p_1-p_3)^2, \label{eq:lec07:Mt}\\
  \Mamp_u &\sim \frac{\lambda^2}{u - m^2}, & u &= (p_1-p_4)^2. \label{eq:lec07:Mu}
\end{align}
The total amplitude is $\Mamp = \Mamp_s + \Mamp_t + \Mamp_u$.

\subsubsection*{Kinematics in the center-of-mass frame}

Choose the $z$-axis along the beam.  The four four-momenta are:
\begin{align}
  p_1 &= (E,\;0,\;0,\;+p_{\rm CM}), \label{eq:lec07:p1}\\
  p_2 &= (E,\;0,\;0,\;-p_{\rm CM}), \label{eq:lec07:p2}\\
  p_3 &= (E,\;p_{\rm CM}\sin\theta,\;0,\;+p_{\rm CM}\cos\theta), \label{eq:lec07:p3}\\
  p_4 &= (E,\;-p_{\rm CM}\sin\theta,\;0,\;-p_{\rm CM}\cos\theta), \label{eq:lec07:p4}
\end{align}
where all particles have equal mass $m$, so $E^2 = p_{\rm CM}^2 + m^2$.

\subsubsection*{Computing the Mandelstam variables}

\textbf{Computing $s$.}  We have $p_1+p_2 = (2E,\bvec{0})$, so:
\begin{equation}
  s = (p_1+p_2)^2 = 4E^2.
  \label{eq:lec07:s-result}
\end{equation}

\textbf{Computing $t$.}  We compute $p_1 - p_3$ component by component:
\[
  p_1 - p_3 = \bigl(0,\; -p_{\rm CM}\sin\theta,\; 0,\; p_{\rm CM}(1-\cos\theta)\bigr).
\]
Then:
\begin{align}
  t = (p_1-p_3)^2
  &= 0^2 - p_{\rm CM}^2\sin^2\theta - p_{\rm CM}^2(1-\cos\theta)^2 \notag\\
  &= -p_{\rm CM}^2\bigl[\sin^2\theta + 1 - 2\cos\theta + \cos^2\theta\bigr] \notag\\
  &= -p_{\rm CM}^2\bigl[2 - 2\cos\theta\bigr] \notag\\
  &= \boxed{-2p_{\rm CM}^2(1-\cos\theta).}
  \label{eq:lec07:t-result}
\end{align}
(Note: the Minkowski metric introduces an overall minus sign on the spatial part; the zero component is zero here since all masses are equal.)

\textbf{Computing $u$.}  Similarly, $p_1 - p_4 = (0,\;p_{\rm CM}\sin\theta,\;0,\;p_{\rm CM}(1+\cos\theta))$, giving:
\begin{equation}
  u = -p_{\rm CM}^2(1+\cos\theta)^2 - p_{\rm CM}^2\sin^2\theta
  = \boxed{-2p_{\rm CM}^2(1+\cos\theta).}
  \label{eq:lec07:u-result}
\end{equation}

\begin{remark}[Sanity check: $s+t+u = 4m^2$]
One can verify $s+t+u = 4E^2 - 2p_{\rm CM}^2(1-\cos\theta) - 2p_{\rm CM}^2(1+\cos\theta) = 4E^2 - 4p_{\rm CM}^2 = 4m^2$, as required by the general identity $s+t+u=\sum_i m_i^2$.
\end{remark}

\subsubsection*{The full differential cross-section}

Substituting into $\Mamp = \Mamp_s + \Mamp_t + \Mamp_u$ and then into \eqref{eq:lec07:dsigma-equal-mass}:
\begin{equation}
  \dcross = \frac{1}{64\pi^2 s}\,\left|
    \frac{\lambda^2}{s-m^2}
    + \frac{\lambda^2}{t-m^2}
    + \frac{\lambda^2}{u-m^2}
  \right|^2.
  \label{eq:lec07:dsigma-full}
\end{equation}
Explicitly, using $s = 4E^2$, $t = -2p_{\rm CM}^2(1-c)$, $u = -2p_{\rm CM}^2(1+c)$ where $c\equiv\cos\theta$:
\begin{equation}
  |\Mamp|^2 = \lambda^4\left(
    \frac{1}{4E^2 - m^2}
    + \frac{1}{-2p_{\rm CM}^2(1-c)-m^2}
    + \frac{1}{-2p_{\rm CM}^2(1+c)-m^2}
  \right)^2.
  \label{eq:lec07:M2-explicit}
\end{equation}
Each channel depends differently on $\cos\theta$: the $s$-channel is flat in angle, while the $t$- and $u$-channels peak in the forward and backward directions respectively.

%%──────────────────────────────────────────────────────────────────
\subsection*{Symmetry Factors via Contraction Counting}

\noindent\textbf{Physical intuition:} When writing down $\Mamp = \Mamp_s + \Mamp_t + \Mamp_u$, one must ensure that the three diagrams enter with the correct relative numerical coefficients.  These come from counting the number of distinct ways to ``wire up'' the external legs to the interaction vertices, using a procedure called \emph{Wick contractions}.

\subsubsection*{The method}

Each $\phi^3$ vertex contributes $\lambda\phi\phi\phi$ — three distinguishable leg ``slots''.  For $\phi\phi\to\phi\phi$ at order $\lambda^2$ there are two vertices, each with three legs, giving $3+3=6$ slots in total.  The four external states ($p_1, p_2, p_3, p_4$) and the one internal propagator must each connect a specific pair of slots.

To count the symmetry factor for a given diagram, proceed as follows:
\begin{enumerate}
  \item For each external leg, count the number of available vertex slots it can connect to (given the topology of the diagram), and write down that number.
  \item For the internal propagator, count the number of ways the two remaining slots can connect to each other (usually $1$, since they are forced by topology once the external legs are placed).
  \item Multiply all the numbers together.
\end{enumerate}

\subsubsection*{$s$-channel: diagram A}

External legs $p_1$ and $p_2$ both connect to vertex~1; $p_3$ and $p_4$ both connect to vertex~2; the internal line connects the two vertices.

\begin{center}
\begin{tabular}{lcc}
  \toprule
  Leg & Vertex & Available slots\\
  \midrule
  $p_1$ & vertex 1 & 3\\
  $p_2$ & vertex 1 & 2 (one slot taken)\\
  $p_3$ & vertex 2 & 3\\
  $p_4$ & vertex 2 & 2 (one slot taken)\\
  internal & v1--v2 & 1 (fully determined)\\
  \bottomrule
\end{tabular}
\end{center}
Total: $3 \times 2 \times 3 \times 2 \times 1 = 36$.

\subsubsection*{$t$-channel: diagram B}

External legs $p_1$ and $p_3$ connect to vertex~1; $p_2$ and $p_4$ connect to vertex~2.  The same counting gives:
\[
  3\times 2\times 3\times 2\times 1 = 36.
\]

\subsubsection*{$u$-channel: diagram C}

External legs $p_1$ and $p_4$ connect to vertex~1; $p_2$ and $p_3$ connect to vertex~2.  Again:
\[
  3\times 2\times 3\times 2\times 1 = 36.
\]

\subsubsection*{Conclusion}

All three diagrams have the same symmetry factor (36), so they enter $\Mamp$ with equal coefficients — no extra numerical prefactors between channels.  This confirms that the sum $\Mamp = \Mamp_s + \Mamp_t + \Mamp_u$ with unit relative weights is correct.

\begin{remark}[Interchange of identical vertices]
One might worry about the factor arising from swapping the two vertices (i.e., relabeling vertex~1 as vertex~2).  In full perturbation theory, every $n$-th order term in the expansion of $e^{iS_{\rm int}}$ comes with a factor $1/n!$ from the Taylor expansion.  For two vertices, this $1/2!$ exactly cancels the 2 ways of labeling the vertices.  Therefore, \textbf{one never needs to include an extra factor for interchanging identical vertices}; it is automatically cancelled by the perturbative expansion.
\end{remark}

\dfn{Wick contraction}{The procedure of assigning external states to vertex legs and connecting remaining legs via internal propagators is called a \textbf{Wick contraction}.  Each distinct complete contraction corresponds to one Feynman diagram, and the number of ways to perform a given contraction gives the numerical (symmetry) factor for that diagram.}

%%──────────────────────────────────────────────────────────────────
\subsection*{Dimensional Analysis of Matrix Elements}

\noindent\textbf{Physical intuition:} Natural units relate mass, energy, time, and length.  A quick dimensional check — before evaluating a single loop integral or performing any computation — can immediately reveal algebra errors.

\subsubsection*{Units of observables}

In natural units ($\hbar = c = 1$):
\begin{itemize}
  \item \textbf{Decay rate} $\DecayRate$: has units of $[{\rm time}]^{-1} = [{\rm mass}]^1$.
  \item \textbf{Cross-section} $\sigma$: has units of area $= [{\rm length}]^2 = [{\rm mass}]^{-2}$.
\end{itemize}

\begin{remark}[The barn]
Cross-sections are sometimes quoted in \textbf{barns}: $1\,{\rm barn} = 10^{-24}\,{\rm cm}^2$.  This is a convenient unit for nuclear and particle physics cross-sections.
\end{remark}

\subsubsection*{Mass dimension of the matrix element}

\thm{Dimension of $\Mamp$}{For a process involving $n$ external particles (initial plus final), regardless of how they are distributed between initial and final states, the matrix element has mass dimension:
\begin{equation}
  [\Mamp] = \text{mass}^{4-n}.
  \label{eq:lec07:Mdim}
\end{equation}}

Examples:
\begin{itemize}
  \item $1\to2+3$ (decay, $n=3$): $[\Mamp] = {\rm mass}^1$.
  \item $1+2\to3+4$ (scattering, $n=4$): $[\Mamp] = {\rm mass}^0$ (dimensionless).
  \item $1\to2+3+4$ (three-body decay, $n=4$): $[\Mamp] = {\rm mass}^0$ (dimensionless).
\end{itemize}

\subsubsection*{Consistency check}

One can verify these dimensions are consistent with the formulas for $\DecayRate$ and $\sigma$:
\begin{itemize}
  \item \textbf{Two-body decay} \eqref{eq:lec07:Gamma2body}: $[\DecayRate] = [p_*]/[M_1^2]\cdot[\Mamp]^2 = {\rm mass}\cdot{\rm mass}^{-2}\cdot{\rm mass}^2 = {\rm mass}^1$.\ \checkmark
  \item \textbf{$2\to2$ cross-section} \eqref{eq:lec07:dsigma-equal-mass}: $[\dcross] = [\Mamp]^2/[s] = {\rm mass}^0/{\rm mass}^2 = {\rm mass}^{-2}$. \checkmark
\end{itemize}

%%──────────────────────────────────────────────────────────────────
\subsection*{Summary: Structure of Perturbative QFT}

The framework we have built can be summarized as follows:

\begin{itemize}
  \item \textbf{Quadratic (free) terms} in $\lagr$: terms with exactly two fields, $\lagr \supset \frac{1}{2}(\partial\phi)^2 - \frac{1}{2}m^2\phi^2$.  These describe \textbf{free fields} — particles that propagate without interacting.  Free particles cannot be created or destroyed; there are no nontrivial scattering or decay processes.

  \item \textbf{Higher-order (interaction) terms}: terms with three or more fields, e.g., $\lagr \supset \lambda\phi^3$ or $\lambda\phi^4$.  These are treated as \textbf{perturbations} of the free theory, valid when the coupling $\lambda$ is small.  They introduce interactions: particles can now be created, destroyed, scattered, or decay.

  \item \textbf{Feynman diagrams}: a systematic graphical tool for computing the transition amplitude (matrix element $\Mamp$) for any process, order by order in $\lambda$.  Each diagram corresponds to a term in perturbation theory.

  \item \textbf{From $\Mamp$ to observables}: the matrix element squared $|\Mamp|^2$, combined with the appropriate phase-space integral, gives physical observables — decay rates $\DecayRate$ and cross-sections $\sigma$ — that can be compared to experiment.

  \item \textbf{Experiment}: once calculations are in hand, one runs experiments to measure $\DecayRate$ or $\sigma$ and compare with theory to test whether the theory describes nature.
\end{itemize}

%%──────────────────────────────────────────────────────────────────
\subsection*{Fields We Care About}

Before diving into fermions, it is useful to catalog the types of fields that appear in the Standard Model and in any relativistic QFT:

\begin{center}
\renewcommand{\arraystretch}{1.4}
\begin{tabular}{llll}
  \toprule
  Spin & Field type & Symbol & Example \\
  \midrule
  $0$ & Scalar & $\phi$ & Higgs boson $H$\\
  $1/2$ & Fermion (Dirac spinor) & $\psi$, $\psi_L$, $\psi_R$ & Quarks, leptons (electrons, neutrinos, etc.)\\
  $1$ & Vector (gauge boson) & $A^\mu$ & Photon $\gamma$, $W^\pm$, $Z^0$, gluons\\
  \bottomrule
\end{tabular}
\end{center}

Scalars and vectors are in some sense analogous in their treatment; fermions require substantially more formalism.  The gauge bosons (spin-1) will be covered when we discuss local (gauge) symmetries later in the course.

%%──────────────────────────────────────────────────────────────────
\subsection*{Introduction to Fermions}

\noindent\textbf{Physical intuition:} For bosons (spin-0, spin-1), we started from the Klein--Gordon equation.  For fermions (spin-$\frac{1}{2}$), the appropriate starting point is the Dirac equation.  Fermions obey anti-commutation relations (rather than commutation relations) for their creation/annihilation operators, and their fields carry spinor indices.

\subsubsection*{The Dirac equation}

The equation of motion for a free spin-$\frac{1}{2}$ field is the \textbf{Dirac equation}:
\begin{equation}
  \boxed{\bigl(i\gamma^\mu\partial_\mu - m\bigr)\psi(x) = 0.}
  \label{eq:lec07:Dirac}
\end{equation}

The field $\psi(x)$ is called a \textbf{Dirac fermion} (or Dirac spinor); it obeys the Dirac equation because of the specific structure of $\gamma^\mu$.

\subsubsection*{The Dirac spinor $\psi$}

$\psi(x)$ is a \textbf{four-component} field:
\begin{equation}
  \psi(x) = \begin{pmatrix}\psi_1(x)\\\psi_2(x)\\\psi_3(x)\\\psi_4(x)\end{pmatrix}.
  \label{eq:lec07:spinor-components}
\end{equation}
\textbf{Important}: $\psi$ is \emph{not} a four-vector.  A four-vector transforms under Lorentz boosts and rotations as $V^\mu \to \Lambda^\mu_{\ \nu}V^\nu$; $\psi$ transforms differently (as a spinor representation of the Lorentz group).  The four components of $\psi$ are internal degrees of freedom — two spin states times two particle/antiparticle states.  $\psi$ is also called a \textbf{Dirac spinor}.

From the quantum-mechanical perspective: fermions are excitations of fields whose mode expansions use \emph{fermionic} harmonic oscillator operators, satisfying anti-commutation relations $\{a, a^\dagger\} = 1$ rather than commutation relations.

\subsubsection*{Gamma matrices (Dirac matrices)}

\noindent\textbf{Physical intuition:} In non-relativistic quantum mechanics, spin-$\frac{1}{2}$ particles are described using Pauli matrices $\sigma^i$.  The Dirac matrices $\gamma^\mu$ are the relativistic generalization: four $4\times4$ matrices (one for each spacetime direction) that encode both the Pauli matrices and the relativistic kinematics.

The index $\mu$ in $\gamma^\mu$ runs from $0$ to $3$, so $\gamma^\mu$ represents four matrices:
\[
  \gamma^0,\; \gamma^1,\; \gamma^2,\; \gamma^3.
\]
Each $\gamma^\mu$ is a $4\times4$ matrix.  In the \textbf{standard (Dirac) basis}, they are written in $2\times2$ block form:
\begin{equation}
  \gamma^\mu = \begin{pmatrix}0 & \sigma^\mu \\ \bar\sigma^\mu & 0\end{pmatrix},
  \label{eq:lec07:gamma-block}
\end{equation}
where
\begin{align}
  \sigma^\mu &\equiv \bigl(\mathbf{1}_{2\times2},\;\sigma^1,\;\sigma^2,\;\sigma^3\bigr), \label{eq:lec07:sigma-mu}\\
  \bar\sigma^\mu &\equiv \bigl(\mathbf{1}_{2\times2},\;-\sigma^1,\;-\sigma^2,\;-\sigma^3\bigr), \label{eq:lec07:sigmabar-mu}
\end{align}
and $\sigma^i$ are the three $2\times2$ Pauli matrices.  Explicitly:
\begin{equation}
  \gamma^0 = \begin{pmatrix}0 & \mathbf{1}\\\mathbf{1} & 0\end{pmatrix}, \qquad
  \gamma^i = \begin{pmatrix}0 & \sigma^i\\-\sigma^i & 0\end{pmatrix},\quad i=1,2,3.
  \label{eq:lec07:gamma-explicit}
\end{equation}

\begin{remark}[Different bases]
The representation \eqref{eq:lec07:gamma-block} is sometimes called the \textbf{Weyl} or \textbf{chiral basis}.  Peskin \& Schroeder use this basis.  Other texts (e.g., Griffiths) use a different basis (the \textbf{Dirac basis}) related by a unitary transformation.  Physical results are basis-independent.
\end{remark}

\subsubsection*{Clifford algebra}

The gamma matrices satisfy the fundamental anti-commutation relation:
\begin{equation}
  \boxed{\anticomm{\gamma^\mu}{\gamma^\nu} \equiv \gamma^\mu\gamma^\nu + \gamma^\nu\gamma^\mu = 2g^{\mu\nu}\,\mathbf{1}_{4\times4},}
  \label{eq:lec07:Clifford}
\end{equation}
where $g^{\mu\nu} = \mathrm{diag}(+1,-1,-1,-1)$ is the Minkowski metric.  This is the \textbf{Clifford algebra}.  It is the defining property of the Dirac matrices — any $4\times4$ matrix representation satisfying \eqref{eq:lec07:Clifford} gives an equivalent formulation of the Dirac equation.

\subsubsection*{The $\gamma^5$ matrix}

An additional important matrix is:
\begin{equation}
  \gamma^5 \equiv i\,\gamma^0\gamma^1\gamma^2\gamma^3.
  \label{eq:lec07:gamma5-def}
\end{equation}
In the standard basis:
\begin{equation}
  \gamma^5 = \begin{pmatrix}-\mathbf{1} & 0 \\ 0 & \mathbf{1}\end{pmatrix}.
  \label{eq:lec07:gamma5-matrix}
\end{equation}
Key properties:
\begin{itemize}
  \item $(\gamma^5)^2 = \mathbf{1}$ (it squares to the identity).
  \item $(\gamma^5)^\dagger = \gamma^5$ (it is Hermitian).
  \item $\anticomm{\gamma^5}{\gamma^\mu} = 0$ for $\mu = 0,1,2,3$ (it anti-commutes with all $\gamma^\mu$).
\end{itemize}
$\gamma^5$ is the chirality operator.  Its eigenvalues $\pm1$ define right-handed ($+$) and left-handed ($-$) chirality states (the Weyl spinors, discussed below).

\subsubsection*{Dirac conjugate $\bar\psi$}

For building Lorentz-invariant bilinears, one defines the \textbf{Dirac conjugate}:
\begin{equation}
  \bar\psi \equiv \psi^\dagger\,\gamma^0.
  \label{eq:lec07:Dirac-conjugate}
\end{equation}
One always works with $\bar\psi$ rather than $\psi^\dagger$; the factor of $\gamma^0$ is needed to ensure Lorentz invariance.  For example, $\bar\psi\psi$ is a Lorentz scalar, while $\psi^\dagger\psi$ is not.

\subsubsection*{Dirac Lagrangian}

The Lagrangian that yields the Dirac equation as its equation of motion is:
\begin{equation}
  \boxed{\lagr_{\rm Dirac} = \bar\psi\,(i\gamma^\mu\partial_\mu - m)\,\psi.}
  \label{eq:lec07:Dirac-Lagrangian}
\end{equation}
Applying the Euler--Lagrange equation with respect to $\bar\psi$ (treated as an independent degree of freedom) gives:
\[
  \frac{\partial\lagr}{\partial\bar\psi} = (i\gamma^\mu\partial_\mu - m)\psi = 0,
\]
which is precisely the Dirac equation \eqref{eq:lec07:Dirac}.  We treat $\psi$ and $\bar\psi$ as independent degrees of freedom, exactly as one treats $\phi$ and $\phi^\dagger$ independently for a complex scalar.

\subsubsection*{Weyl spinors}

The Dirac spinor $\psi$ decomposes into two \textbf{two-component Weyl spinors}:
\begin{equation}
  \psi = \begin{pmatrix}\psi_L \\ \psi_R\end{pmatrix},
  \label{eq:lec07:Weyl-decomp}
\end{equation}
where $\psi_L$ (the upper two components) is called the \textbf{left-handed Weyl spinor} and $\psi_R$ (the lower two components) is the \textbf{right-handed Weyl spinor}.  Each is a two-component object.

\dfn{Weyl spinors}{The left- and right-handed Weyl spinors $\psi_{L,R}$ are the eigenstates of $\gamma^5$:
\begin{equation}
  \gamma^5\psi_L = -\psi_L, \qquad \gamma^5\psi_R = +\psi_R.
  \label{eq:lec07:chirality-eigenvalues}
\end{equation}
They are the irreducible representations of the Lorentz group for spin-$\frac{1}{2}$ fields.  Under Lorentz transformations, $\psi_L$ and $\psi_R$ transform into themselves (not into each other), making them the basic building blocks of any spinor field.}

\begin{remark}[Why Weyl spinors are fundamental]
Under Lorentz boosts and rotations, $\psi_L$ and $\psi_R$ transform separately (though with related parameters).  The Dirac spinor $\psi$ is therefore a reducible representation: it is the direct sum of a left-handed and a right-handed Weyl representation.  In the Standard Model, it turns out that the matter fields (quarks and leptons) have different quantum numbers for $\psi_L$ and $\psi_R$ — left and right handedness are treated differently by the weak force — so the Weyl spinors are the physically natural variables.
\end{remark}

%%──────────────────────────────────────────────────────────────────
\subsection*{To Remember}

\begin{itemize}
  \item \textbf{Two-body decay rate}:
    $\displaystyle\DecayRate_{1\to23} = \frac{p_*}{8\pi M_1^2}\,|\Mamp|^2,$
    where $p_*$ is the CM momentum of either daughter, fixed by the three masses.

  \item \textbf{General cross-section prefactor}: $1/(2E_1\,2E_2\,|\bvec{v}_1-\bvec{v}_2|)$; its Lorentz-invariant form is $1/(4\sqrt{(p_1\cdot p_2)^2-m_1^2m_2^2})$.

  \item \textbf{Equal-mass $2\to2$ cross-section}: $d\sigma/d\Omega = |\Mamp|^2/(64\pi^2 s)$.

  \item \textbf{Mandelstam variables in CM frame} (equal mass $m$): $s = 4E^2$, $t = -2p_{\rm CM}^2(1-\cos\theta)$, $u = -2p_{\rm CM}^2(1+\cos\theta)$.

  \item \textbf{Symmetry factors via contraction counting}: count the number of ways to wire external legs to vertex slots; equal factors for all channels means equal relative weights.

  \item \textbf{Dimension of $\Mamp$}: $[\Mamp] = {\rm mass}^{4-n}$ for an $n$-particle process.

  \item \textbf{Dirac equation}: $(i\Dirac{\partial}-m)\psi = 0$.

  \item \textbf{$\psi$ is a four-component Dirac spinor}, not a four-vector.

  \item \textbf{Gamma matrices} satisfy the Clifford algebra: $\anticomm{\gamma^\mu}{\gamma^\nu} = 2g^{\mu\nu}$.

  \item \textbf{$\gamma^5 = i\gamma^0\gamma^1\gamma^2\gamma^3$}: squares to $+1$, anti-commutes with $\gamma^\mu$, eigenvalues $\pm1$ define chirality.

  \item \textbf{Dirac conjugate}: $\bar\psi = \psi^\dagger\gamma^0$.  Always use $\bar\psi$, not $\psi^\dagger$.

  \item \textbf{Dirac Lagrangian}: $\lagr = \bar\psi(i\Dirac\partial - m)\psi$.

  \item \textbf{Weyl decomposition}: $\psi = (\psi_L,\,\psi_R)^T$; $\psi_L$ and $\psi_R$ are the irreducible Lorentz representations, each two-component.
\end{itemize}

%%──────────────────────────────────────────────────────────────────
\subsection*{Recommended Reading}

\begin{itemize}
  \item \textbf{Griffiths, Chapter~6.3--6.4} — Two-body and three-body decay rates in the toy scalar theory; cross-section formula.  The worked kinematics are directly parallel to this lecture.

  \item \textbf{Peskin \& Schroeder, Chapter~3} — The Dirac equation, Dirac matrices, Clifford algebra, Weyl spinors, and the Dirac Lagrangian.  The lecturer explicitly recommended Chapter~3 for the fermion material in this and subsequent lectures.

  \item \textbf{Halzen \& Martin, Chapter~4--5} — Phase-space kinematics, two-body scattering and decay, followed by the introduction of spinors and Dirac fermions.  A very accessible treatment.

  \item \textbf{Tong, \textit{Lectures on Particle Physics}, Chapter~3}
    (\href{https://www.damtp.cam.ac.uk/user/tong/pp.html}{\texttt{damtp.cam.ac.uk/user/tong/pp.html}}):
    phase space and cross-sections with the same conventions used in this course; Chapter~4 covers the Dirac equation.

  \item \textbf{PDG Review of Kinematics}
    (\href{https://pdg.lbl.gov/2023/reviews/rpp2023-rev-kinematics.pdf}{\texttt{pdg.lbl.gov/2023/reviews/rpp2023-rev-kinematics.pdf}}):
    compact reference for two-body decay kinematics, Mandelstam variables, and cross-section formulas.
\end{itemize}

\end{document}
